% !TeX root = ../main.tex

\chapter*{摘要}\addcontentsline{toc}{chapter}{摘要}\label{ux6458ux8981}

当前，随着人工智能的蓬勃发展，大型语言模型已在代码生成、错误诊断、自然语言交互及个性化反馈等多个维度展现出显著潜力，这为中小学生编程教育开辟了新的技术探索方向。传统的编程教学体系在资源获取的普适性、个性化辅导的及时性、学习进度的量化追踪以及即时反馈的效率等环节存在固有的局限，这容易在学习的初阶阶段造成学生兴趣减退乃至学习挫败感。为应对上述挑战，本文构建并实现了一套面向中小学生的智能编程学习平台。

平台采纳了前后端分离架构设计：前端基于 Nuxt 3、Vue 3、Tailwind CSS 和 DaisyUI 实现用户交互界面，后端则依托 FastAPI 与 SQLAlchemy 构建核心业务接口与数据管理。数据持久层采用了 SQL Server，并通过 Redis 实现缓存加速，同时利用阿里云 OSS 进行教学资源的存储。功能模块涵盖了用户身份认证、课程知识点编排、在线编程实践、任务管理、师生即时互动、教学资源管理等基本组件。更重要的是，平台深度集成了大模型能力，以支撑自由问答、代码逻辑拆解、编译错误分析、问题分解启发、抽象建模和算法设计层级指导等智能化辅学服务。鉴于中小学生在认知发展上的特定特点，系统在提示词的策略设计、输出语气的校准和错误反馈的模式选择上进行了精细化的调整。

本文围绕从系统需求分析到总体设计、再到核心业务模块的详细实现与功能验证展开论述。从工程落地角度看，该平台已能为学生的日常编程练习及教师的教学管理工作提供一个相对完整的操作支撑。它在代码的稳定运行、错误的精准反馈、学习进度的可量化跟踪以及资源的有效组织等方面奠定了坚实的实践基础，同时，也为后续引入基于 RAG 的私域知识增强、认知模板评估机制以及更精细的个性化推荐系统预留了合理的扩展接口。

\textbf{关键词}：大语言模型；编程教育；智能辅学；在线编程；教学平台

\chapter*{ABSTRACT}\addcontentsline{toc}{chapter}{ABSTRACT}\label{abstract}

The swift advancement of Artificial Intelligence, particularly Large Language Models (LLMs), has yielded notable capabilities across domains such as code generation, diagnostic error spotting, natural language interaction, and personalized tutoring, effectively charting a novel technical route for programming education among primary and secondary school students. Conventional programming pedagogy often encounters structural limitations regarding the accessibility of learning materials, the immediacy of individualized mentorship, the tracking of learning progression, and the efficiency of real-time feedback. These deficiencies frequently lead to decreased student engagement or feelings of learning discouragement during initial programming phases. Addressing these gaps, this work designs and implements an intelligent programming learning platform specialized for K-12 learners, leveraging LLMs.

The platform operates on a separated front-end and back-end architecture. The front end utilizes Nuxt 3, Vue 3, Tailwind CSS, and DaisyUI to manage the user interface, while the back end relies on FastAPI and SQLAlchemy to govern business APIs and data lifecycle. Data storage is handled by SQL Server, augmented with Redis for caching efficiency and Alibaba Cloud OSS for static resource hosting. Key functionalities include user authentication, curriculum structuring, online coding practice, task tracking, real-time communication channels, and resource administration. Crucially, the system integrates LLM services to power intelligent tutoring, covering functions like open-ended Q\&A, code logic deconstruction, compile error analysis, problem decomposition guidance, abstract modeling support, and algorithm structuring advice. Given the specific cognitive profiles of young learners, the system incorporates targeted fine-tuning in prompt engineering, output tone control, and error feedback methodology.

This thesis details the entire process, from initial requirements analysis and overall architecture design to the detailed implementation and functional verification of core modules. In terms of engineering viability, the platform provides a robust operational backbone supporting students' daily coding practice and teachers' administrative needs, with solid foundations in code execution stability, accurate error reporting, quantitative progress monitoring, and resource organization. Furthermore, the architecture is designed to seamlessly accommodate future enhancements, such as RAG-based private knowledge augmentation, cognitive template evaluation, and advanced personalized recommendation engines.

\textbf{Key words}: Large Language Model; Programming Education; Intelligent Tutoring; Online Programming; Learning Platform

\chapter{引言}\label{ux5f15ux8a00}

\section{研究背景}\label{ux7814ux7a76ux80ccux666f}

在教育数字化持续推进的背景下，编程教育逐渐成为中小学生信息素养培养的重要组成部分。通过编程学习，学生不仅能够掌握基础语法知识，还能在实践过程中培养逻辑思维、抽象建模能力和问题求解能力。然而，中小学生在编程学习初期通常面临语法理解困难、报错定位能力不足、学习反馈不及时和教师个性化指导成本高等问题，这些问题会直接影响学习体验和持续学习意愿。

近年来，大语言模型在自然语言理解与生成、代码理解与生成、交互式问答等方面表现突出，为构建智能化编程学习平台提供了现实条件。已有教育综述研究指出，LLM 正在从学生辅助、教师支持、自适应学习和教育工具等多个方向重塑教育应用形态{[}1{]}。在计算机科学教育领域，系统综述也表明，LLM 已逐步被用于代码生成、代码解释、调试辅助与个性化学习支持{[}2{]}。将大模型应用于编程教育场景，可以在一定程度上承担代码解释、错误诊断、思路引导与启发式提问等功能，从而缓解传统教学中反馈滞后和个别辅导不足的问题。因此，研究并实现一套面向中小学生的智能编程学习平台，具有较强的现实意义。

\section{研究意义}\label{ux7814ux7a76ux610fux4e49}

本文的研究意义主要体现在以下几个方面。

首先，在研究层面，现有关于 LLM 赋能编程教育的成果，较多聚焦于高等教育、单一教学任务或工具效果评估，对于面向中小学生、覆盖课程学习、在线编程、错误诊断、教师干预和过程记录完整链条的平台化研究仍相对不足{[}2{]}{[}3{]}{[}12{]}。因此，本文从系统实现与教学适配两个维度展开研究，能够为中小学生智能编程学习平台的设计提供较为具体的实现性参考。

其次，在教学层面，中小学生在编程入门阶段普遍存在概念抽象、错误定位困难和反馈滞后等问题。Tang 等的研究表明，在教学设计合理的前提下，LLM 支持能够提升学生的编程表现、学习兴趣与编程自我效能感{[}4{]}。据此，本人认为，将大模型能力嵌入学生日常学习流程，而不是将其作为孤立的问答工具，更有可能形成连续、稳定的学习支持机制。

再次，在技术层面，本文将 Nuxt 3、FastAPI、SQLAlchemy、SQL Server、Redis 与大模型接口进行整合，构建了面向教育场景的前后端协同系统。该实现过程不仅涉及常规业务模块建设，还涉及代码安全执行、提示约束、行为留痕与角色权限控制等问题，因此具有一定的工程实践价值。

最后，在应用层面，LLM 虽然能够提升反馈速度，但同时也带来输出不稳定、学习依赖和责任边界模糊等风险{[}8{]}{[}9{]}{[}10{]}{[}11{]}。本人在本文中并未将大模型定位为教师替代者，而是将其作为教师可调控的辅学组件加以设计。这种定位更符合中小学生教学场景对安全性、可控性和可追踪性的基本要求，也使平台具备更现实的推广价值和扩展空间。

\section{国内外研究现状}\label{ux56fdux5185ux5916ux7814ux7a76ux73b0ux72b6}

国内外关于大模型与编程教育结合的研究主要集中在以下几个方向。

\begin{enumerate}
\def\labelenumi{(\arabic{enumi})}
\tightlist
\item
  大模型在代码生成、代码解释和错误诊断中的应用研究。
\item
  K-12 编程教育平台的教学支持机制研究。
\item
  个性化学习反馈、学习行为分析与教学干预研究。
\item
  教育场景中大模型使用的可信性、安全性与伦理问题研究。
\end{enumerate}

从综述类研究来看，Wang 等对教育领域的 LLM 应用进行了系统总结，指出其在学生支持、教师辅助和自适应学习方面具有广阔潜力，但也伴随可靠性、偏差和伦理风险{[}1{]}。Raihan 等针对计算机科学教育开展系统文献综述，发现 LLM 已经广泛参与编程学习、代码理解和课程支持，但对实际学习成效与教学适配性的研究仍在发展之中{[}2{]}。Pitts 等进一步聚焦编程教育场景，认为 LLM 在形成性反馈、自动评估与知识建模方面具有应用价值，但真正有效的设计往往依赖教师参与和人在回路机制{[}3{]}。

从 K-12 与中小学编程教育的实证研究来看，Tang 等在中国初中场景中的研究表明，LLM 支持能够提升学生的编程表现、学习兴趣和编程自我效能感{[}4{]}。Chen 等提出的 ChatScratch 系统将 AI 增强机制与可视化编程环境结合，显示出儿童在自主学习和作品完成质量方面的积极变化{[}5{]}。Ma 等则从调试教学角度提出，LLM 更适合作为协同式、可教学的辅助对象，通过引导学生提出假设、定位错误和解释修改理由来促进调试能力发展，而不是直接替代学生完成问题求解{[}6{]}。

与此同时，现有研究也揭示出不少限制。Gonzalez-Maldonado 等在 K-12 块编程教学中发现，仅依靠提示工程并不足以让 LLM 稳定完成项目生成和自动评分任务，模型在语法和结果可靠性上仍存在明显缺陷{[}7{]}。Pirzado 等的系统综述指出，LLM 在代码理解、调试支持、学习依赖和反馈可信性等方面仍存在较多潜在风险{[}8{]}。Humble 进一步从计算教育治理视角指出，生成式人工智能的教育应用必须同时考虑优势与威胁，不能仅以效率提升作为唯一依据{[}9{]}。此外，Jeon 与 Lee 强调，教育场景中更现实的路径是构建教师与模型的互补关系，而非将模型简单替代为自动教师{[}11{]}。因此，虽然大模型在降低编程学习门槛、提升即时反馈能力方面具有明显优势，但在中小学生教育场景中，仍面临回答准确性波动、输出难度不稳定、教学语气控制不足以及学习过程数据与资源整合不充分等问题。基于此，构建一个既具备基础教学管理能力、又能融入大模型智能辅学能力的平台，仍具有较强的研究必要性。

\section{研究内容与目标}\label{ux7814ux7a76ux5185ux5bb9ux4e0eux76eeux6807}

本文以``基于大模型的中小学生智能编程学习平台设计与实现''为研究主题，围绕以下内容展开。

\begin{enumerate}
\def\labelenumi{(\arabic{enumii})}
\tightlist
\item
  设计并实现面向学生与教师的编程学习平台。
\item
  构建课程、知识点、任务、资源和行为记录等核心业务模块。
\item
  集成大模型能力，实现代码解释、报错分析、问题分解、抽象建模和算法设计引导。
\item
  实现在线代码运行与安全控制机制。
\item
  对系统功能与运行效果进行测试分析。
\end{enumerate}

\section{论文结构安排}\label{ux8bbaux6587ux7ed3ux6784ux5b89ux6392}

本文共分为七章。第一章为引言，介绍研究背景、研究意义、研究现状以及论文结构。第二章介绍系统设计与实现所涉及的关键技术。第三章从用户角色、功能需求和非功能需求等方面对系统进行需求分析。第四章给出系统总体设计，包括总体架构、模块设计和数据库设计。第五章介绍系统各模块的详细实现。第六章对系统进行功能测试与结果分析。第七章总结全文工作，并对后续改进方向进行展望。

\chapter{相关技术与理论基础}\label{ux76f8ux5173ux6280ux672fux4e0eux7406ux8bbaux57faux7840}

\section{大语言模型与智能辅学}\label{ux5927ux8bedux8a00ux6a21ux578bux4e0eux667aux80fdux8f85ux5b66}

大语言模型具备较强的自然语言理解与生成能力，在教育场景中可承担问答辅助、内容解释、错误分析和学习引导等任务{[}1{]}{[}10{]}。对于编程教育而言，大模型不仅可以直接回答编程相关问题，还可以将复杂概念转化为更适合初学者理解的表达方式，从而在一定程度上缓解学生对专业术语和抽象知识的理解困难。已有研究表明，LLM 在编程教育中的典型应用包括代码解释、调试支持、形成性反馈、自动评估和知识建模{[}2{]}{[}3{]}。

但需要指出的是，LLM 的教学价值并不等同于其可以直接替代教师。相关研究已经表明，模型在真实性、稳定性、可追责性与教学适配性方面仍存在明显边界{[}7{]}{[}8{]}{[}9{]}{[}10{]}。尤其是在编程学习场景中，模型可能生成表面连贯但实质有误的解释，也可能过早给出完整答案，削弱学生自行分析、调试和抽象概括的过程{[}6{]}{[}8{]}。因此，若缺少明确的教学约束，LLM 很容易从``辅助学习工具''滑向``替代思考工具''。

基于上述认识，本文在平台设计中并未将 LLM 设定为自动判分器或唯一答案源，而是将其定位为过程性辅学模块。本人更强调其支架作用，即在学生卡点处提供适量提示，在错误定位处提供结构化解释，在问题分解处提供思路引导，并在教师管理端保留审阅、追踪与干预空间{[}3{]}{[}6{]}{[}11{]}。这种定位更符合中小学生编程学习的认知发展特点，也更有利于维持课堂教学中的教师主导性。

\section{前端开发技术}\label{ux524dux7aefux5f00ux53d1ux6280ux672f}

本系统前端采用 Nuxt 3 作为核心框架，基于 Vue 3 组合式 API 完成页面逻辑编写，使用 Tailwind CSS 与 DaisyUI 实现界面设计，并通过 TypeScript 提升代码可维护性和类型安全性。该技术组合既能满足交互式页面开发需求，也具备较好的组件化组织能力。

\section{后端开发技术}\label{ux540eux7aefux5f00ux53d1ux6280ux672f}

系统后端采用 FastAPI 作为核心 Web 框架，具备接口开发效率高、性能较好和文档生成方便等特点。数据库访问通过 SQLAlchemy 实现 ORM 映射，业务认证采用 JWT 机制完成身份校验。同时，系统使用 Redis 进行部分缓存处理，以降低高频读取场景的响应压力。

\section{数据存储与资源管理技术}\label{ux6570ux636eux5b58ux50a8ux4e0eux8d44ux6e90ux7ba1ux7406ux6280ux672f}

平台结构化数据主要存储在 SQL Server 中，包括用户、课程、任务、知识点、资源、评论消息、点赞关系与学习记录等信息。对于教学视频等非结构化资源，系统结合阿里云 OSS 进行对象存储，以提高资源访问与管理效率。

\section{在线编程与安全执行技术}\label{ux5728ux7ebfux7f16ux7a0bux4e0eux5b89ux5168ux6267ux884cux6280ux672f}

平台实现了在线代码运行功能。为了保证系统稳定性与使用安全，后端在代码执行前会先进行语法检查和安全检查，过滤危险导入与危险函数调用，并通过超时控制和执行次数限制降低恶意代码或死循环带来的风险。

\chapter{系统需求分析}\label{ux7cfbux7edfux9700ux6c42ux5206ux6790}

\section{应用场景分析}\label{ux5e94ux7528ux573aux666fux5206ux6790}

本平台主要面向中小学生编程学习场景，兼顾教师教学管理需求。其典型应用场景包括课内编程教学、课后作业练习、在线答疑、资源分享和学习过程跟踪等。

从实际教学流程看，本平台的使用并非集中在单一时点，而是贯穿``课前准备---课中讲解---课后巩固---过程评价''多个环节。课前，教师可通过平台整理课程资源、配置任务要求和关联知识点；课中，学生可在统一环境中完成代码编写、运行调试和即时提问；课后，教师则能够依据学习记录查看学生的常见错误与完成进度，并据此调整后续教学安排。换言之，平台承担的并不是单纯的代码运行工具角色，而是一个贯通教学流程的学习支持载体。

在学生侧，本平台尤其适用于编程入门阶段的高频卡点场景。例如，学生面对报错信息时，往往能够看见结果，却难以理解错误含义，更难以判断修改顺序；学生面对综合任务时，也常出现``知道题目要做什么，但不知道第一步从哪里开始''的情况。本人认为，这两类问题恰恰是传统编程教学中最容易造成挫败感的环节，也是引入智能辅学能力最有现实价值的环节。

在教师侧，平台更强调过程管理而非结果收集。与只看最终提交结果的方式相比，教师更需要知道学生在哪个知识点反复出错、在哪个阶段停留时间过长，以及哪些资源真正被学生使用。基于此，平台在设计时不仅关注作业提交与课程展示，也重视消息互动、行为记录和错误追踪等过程性数据。

\section{用户角色分析}\label{ux7528ux6237ux89d2ux8272ux5206ux6790}

系统主要包含以下两类核心用户。

\begin{enumerate}
\def\labelenumi{(\arabic{enumii})}
\tightlist
\item
  学生用户：完成课程学习、代码编写、任务提交、错误分析、消息互动和学习记录查看。
\item
  教师用户：查看学生学习进度、统计错误信息、管理教学资源、发布或维护课程任务并进行师生互动。
\end{enumerate}

上述两类角色在平台中的关注点并不相同。学生用户更关心操作流程是否清晰、反馈是否及时、解释是否易懂，因而其需求主要集中在界面易用性、反馈即时性和辅学内容的可理解性方面。尤其对于编程基础较弱的学习者而言，若系统返回的提示过于抽象，或一次给出过多信息，反而会增加理解负担。

教师用户则更重视平台的组织能力和可管理性。教师不仅需要发布课程、维护知识点和上传资源，还需要借助平台掌握学生学习动态，并对常见问题进行集中处理。因此，教师端不应只是学生端功能的简单镜像，而应当具备更强的管理属性、统计属性和教学干预属性。

从权限控制角度看，学生与教师的角色边界也必须明确。学生侧应以学习和提交为主，教师侧则拥有课程维护、资源管理、学习记录查看等更高权限。只有在角色职责清晰的前提下，平台才能在保证数据安全的同时维持较好的交互秩序。

\section{功能需求分析}\label{ux529fux80fdux9700ux6c42ux5206ux6790}

\subsubsection{3.3.1 学生端功能需求}\label{ux5b66ux751fux7aefux529fux80fdux9700ux6c42}

学生端应满足以下功能需求。

\begin{enumerate}
\def\labelenumi{(\arabic{enumii})}
\tightlist
\item
  登录认证与身份识别。
\item
  查看课程与学习任务。
\item
  在线编写并运行代码。
\item
  获取代码解释与报错分析结果。
\item
  使用智能问答功能寻求帮助。
\item
  与教师或同学进行消息互动。
\item
  查看学习进度、错误记录和推荐资源。
\end{enumerate}

上述功能需求可以进一步归纳为``学什么、怎么练、卡住时怎么办、学完后怎么看''四个层面。首先，学生需要通过课程页和知识点结构明确当前学习对象，即知道自己正在学习哪一部分内容、对应哪些任务以及应达到何种目标。若课程结构展示不清，学生很容易在任务开始前就产生方向性困惑。

其次，在练习过程中，学生需要稳定的在线编程环境。该环境不仅要支持代码输入与运行，还要保证输出结果和错误信息能够及时返回。对于初学者而言，代码运行反馈越延迟，学习链条就越容易中断。因此，在线编程能力并非平台的附属模块，而是学生端体验的核心支点。

再次，当学生遇到理解障碍、运行错误或任务难以拆解时，平台应当提供适量而非过量的辅助。本人认为，学生端智能辅学功能的关键不在于``回答得多完整''，而在于``回答能否让学生接得住''。这意味着系统在设计时应优先支持解释、提示和引导，而不是直接输出大段答案。

最后，学生还需要能够回看自己的学习过程。学习进度、错误记录和历史互动内容不仅有助于学生形成反思，也能在一定程度上减少重复性提问，提高后续学习效率。对于编程初学者而言，看见自己从频繁报错到逐步修正的过程，本身就是一种可感知的成长反馈。

\subsubsection{3.3.2 教师端功能需求}\label{ux6559ux5e08ux7aefux529fux80fdux9700ux6c42}

教师端应满足以下功能需求。

\begin{enumerate}
\def\labelenumi{(\arabic{enumii})}
\tightlist
\item
  查看班级学生学习进度。
\item
  统计学生错误类型与学习表现。
\item
  管理教学资源。
\item
  维护课程任务与知识点内容。
\item
  与学生进行消息交流。
\end{enumerate}

教师端的功能需求可以概括为``可组织、可观察、可干预''三个方面。所谓可组织，是指教师应能够围绕课程主题建立较清晰的内容结构，包括课程分类、知识点关联、任务配置与资源上传，从而使平台内容能够与实际教学进度保持一致。

所谓可观察，是指教师应能够方便地查看学生在学习过程中的关键状态信息，例如任务完成进度、错误分布、活跃情况和典型问题。若教师只能看到最终提交结果，而无法掌握过程数据，那么平台对于教学改进的支撑作用就会被大幅削弱。

所谓可干预，是指教师在发现学生存在集中性问题时，应能够通过消息互动、资源补充或任务调整等方式及时介入。本人认为，面向中小学生的编程学习平台不能只强调学生自学，还必须保留教师干预通道。只有这样，平台中的智能功能才不会与课堂教学脱节，而会成为教师教学能力的延伸。

\subsubsection{3.3.3 智能辅学功能需求}\label{ux667aux80fdux8f85ux5b66ux529fux80fdux9700ux6c42}

系统需要实现面向中小学生的智能辅学能力，包括以下内容。

\begin{enumerate}
\def\labelenumi{(\arabic{enumii})}
\tightlist
\item
  编程相关自由问答。
\item
  代码整体与逐行解释。
\item
  错误信息分析与定位。
\item
  编程任务问题分解。
\item
  抽象建模引导。
\item
  算法设计流程引导。
\end{enumerate}

与传统问答系统相比，本平台的智能辅学需求具有更强的教学场景约束。其一，回答对象是中小学生，因此输出内容必须兼顾正确性、可理解性与适龄性；其二，问题上下文往往与课程任务、学生代码和报错结果绑定，系统不能脱离具体学习情境空泛作答；其三，平台需要避免学生对模型产生过度依赖，因此在部分场景下应优先给出分步提示而非终局答案{[}8{]}{[}11{]}。

基于这一认识，智能辅学功能至少应满足三项要求。第一，能够识别问题类型，并根据不同任务调用不同的回答模板。第二，能够结合学生输入内容生成有针对性的反馈，而不是返回通用化说明。第三，能够对输出进行必要约束，使结果既服务学习，又不突破课堂任务边界。本人认为，这三项要求共同决定了智能辅学模块不能只停留在接口接入层面，而必须在系统设计层面进行专门建模。

\section{非功能需求分析}\label{ux975eux529fux80fdux9700ux6c42ux5206ux6790}

\subsubsection{3.4.1 性能需求}\label{ux6027ux80fdux9700ux6c42}

系统应具备较好的接口响应能力和并发处理能力，能够满足日常教学场景中的多用户访问需求。

考虑到平台的主要使用场景集中在课堂或作业提交时段，系统在性能上不仅要满足普通页面访问，还要关注高频操作的集中出现。例如，某一时间段内多名学生同时打开课程页面、运行代码或请求报错解析，都会对接口响应和任务处理能力提出更高要求。因此，性能需求应重点落在课程访问、代码运行和智能问答等高频路径上，而不是平均分布到所有模块。

\subsubsection{3.4.2 安全需求}\label{ux5b89ux5168ux9700ux6c42}

系统需要保证用户身份认证安全、代码执行安全和教学数据安全，避免恶意访问、非法调用和危险代码执行。

其中，代码执行安全是本平台区别于一般教学信息系统的重要非功能需求。由于平台允许用户提交并运行代码，因此系统必须在执行前进行语法检查、危险调用过滤和超时控制，并在执行后隔离结果回传范围。本人认为，在线编程平台的安全问题更像在实验室中管理化学试剂，重点不在于完全禁止操作，而在于限定边界、控制剂量并设置防护措施。

\subsubsection{3.4.3 可扩展性需求}\label{ux53efux6269ux5c55ux6027ux9700ux6c42}

系统应为后续接入更多课程资源、更丰富的大模型能力以及知识库检索增强等功能保留扩展空间。

从平台演进角度看，可扩展性至少体现在三个方面。第一，业务模块应能够继续增加新的课程类型、任务类型和行为记录类型，而不破坏既有结构。第二，智能辅学模块应能够替换或增加新的模型服务，而不影响整体调用流程。第三，平台应能够为后续接入知识库检索增强、学习者画像和个性化推荐等功能预留接口与数据基础。

\subsubsection{3.4.4 易用性需求}\label{ux6613ux7528ux6027ux9700ux6c42}

由于平台面向中小学生，系统界面和交互流程应尽量简洁直观，智能反馈内容也应贴近学生认知水平。

易用性需求不仅体现为页面是否美观，更体现为操作是否顺手、提示是否明确以及反馈是否容易理解。对于中小学生而言，过多的操作层级、过长的文字说明和过于专业化的表达都会直接提高使用门槛。因此，本平台在交互设计上应尽量缩短核心路径，在反馈设计上应尽量减少术语堆叠，使学生能够把主要注意力放在学习任务本身，而不是放在理解系统如何使用上。

\section{可行性分析}\label{ux53efux884cux6027ux5206ux6790}

\subsubsection{3.5.1 技术可行性}\label{ux6280ux672fux53efux884cux6027}

当前前后端框架、大模型接口能力和数据库技术已经较为成熟，能够支撑本系统的实现。

从开发实现角度看，Nuxt 3、FastAPI、SQLAlchemy、SQL Server 与 Redis 均属于生态较成熟、文档较完善的技术方案，能够覆盖页面交互、接口开发、数据持久化与缓存优化等主要需求。同时，大模型接口的接入方式也已经标准化，便于将自然语言处理能力融入既有业务流程。因此，本文所提出的平台方案在技术栈组合上不存在明显不可实现障碍。

\subsubsection{3.5.2 经济可行性}\label{ux7ecfux6d4eux53efux884cux6027}

系统使用的主要开发工具与框架具有较好的可获得性，部署成本相对可控，具备实际开发基础。

此外，平台的基础模块主要依托通用开发框架实现，软件层面的初始建设成本相对可控。虽然大模型调用会带来一定服务成本，但通过场景限制、请求约束和缓存设计，可以在一定程度上控制无效调用和重复调用。对于毕业设计层面的系统实现而言，现有资源条件能够支撑平台原型的开发、联调与展示。

\subsubsection{3.5.3 运行可行性}\label{ux8fd0ux884cux53efux884cux6027}

系统功能与教学流程匹配度较高，学生与教师均可在已有使用习惯基础上完成操作，具备较好的实际运行可行性。

更重要的是，平台的使用方式与学校日常教学流程具有较高一致性。教师侧的课程发布、任务布置和资源分享，本身就是教学中的常规动作；学生侧的学习、练习、提问和查看反馈，也与日常编程学习行为高度重合。换言之，本平台并未试图重造一套脱离教学实际的流程，而是在原有教学链条上增加了更强的数字化支持，因此具备较好的运行基础。

\chapter{系统总体设计}\label{ux7cfbux7edfux603bux4f53ux8bbeux8ba1}

\section{系统设计目标}\label{ux7cfbux7edfux8bbeux8ba1ux76eeux6807}

系统总体设计遵循模块化、可扩展和安全可控的原则。平台既要满足课程学习、在线编程和师生互动等基础教学需求，又要通过大模型增强反馈能力，从而提升学生在学习过程中的即时性与针对性支持水平。

\section{系统总体架构设计}\label{ux7cfbux7edfux603bux4f53ux67b6ux6784ux8bbeux8ba1}

系统整体采用前后端分离架构，前端负责页面展示与交互逻辑，后端负责业务处理、认证鉴权、代码执行、数据持久化与智能能力调度。系统总体架构如图4-1所示。

\tikzSystemArchitectureFigure

该架构中，前端通过统一接口调用后端服务。后端根据业务将请求分发到用户管理、课程管理、任务管理、消息互动、资源管理、行为记录和智能辅学等模块。对于智能辅学模块，一部分请求通过大模型接口处理自然语言和代码解释任务，另一部分请求则通过安全代码执行器完成在线代码运行与错误返回。

\section{系统功能模块设计}\label{ux7cfbux7edfux529fux80fdux6a21ux5757ux8bbeux8ba1}

根据业务需求，平台功能可以划分为以下几个模块，如图4-2所示。

\tikzFunctionModulesFigure

\section{前后端结构设计}\label{ux524dux540eux7aefux7ed3ux6784ux8bbeux8ba1}

\subsubsection{4.4.1 前端结构设计}\label{ux524dux7aefux7ed3ux6784ux8bbeux8ba1}

前端采用页面与组件结合的方式组织结构，主要页面包括登录页、学生主页、教师主页、课程更新页、事件记录页和事件追踪页。组件层则包含课程分类选择器、课程搜索选择器、知识点选择器、聊天组件和富文本编辑器等。

在前端结构组织上，本文采用``页面承载业务、组件复用交互''的设计思路。页面层主要负责业务流程编排与数据请求触发，组件层则负责可复用的交互单元与局部状态管理。这样设计的好处在于，既能够保证不同业务页面之间的职责边界清晰，又能够在课程筛选、聊天交互、资源编辑等高频场景中减少重复开发工作。

对于学生端页面，设计重点在于缩短学习路径，即尽量减少学生完成``查看课程---进入任务---编写代码---查看反馈''这一主流程时的页面跳转与操作负担。对于教师端页面，设计重点则转向内容维护效率与信息聚合能力，需要在同一界面中尽可能呈现课程编辑、任务配置、资源管理和学习记录查看等常用操作入口。

\subsubsection{4.4.2 后端结构设计}\label{ux540eux7aefux7ed3ux6784ux8bbeux8ba1}

后端以路由模块组织业务接口，并通过统一认证依赖控制访问权限。公开接口主要负责登录与基础班级信息访问，受保护接口则覆盖课程、任务、记录、资源、消息和点赞等模块。该结构有助于实现路由职责清晰、权限边界明确和后续模块扩展。

在后端结构上，系统以业务域划分接口模块，使用户管理、课程管理、任务管理、资源管理、消息互动和学习记录等功能能够在相对独立的边界内演进。本人认为，这种组织方式优于按页面或按角色分散定义接口，因为它更符合服务端对数据和业务规则统一管理的需求，也有利于后续增加新模块时保持接口体系稳定。

同时，后端结构设计中还需要兼顾``通用能力集中化''和``业务能力解耦化''。认证校验、异常处理、统一返回格式等能力应尽量沉淀为公共机制，以避免多个接口重复实现；课程维护、代码运行、智能辅学等业务逻辑则应按模块分别组织，以降低彼此耦合。通过这种方式，系统既能保持整体一致性，又能保留业务扩展弹性。

\section{数据库设计}\label{ux6570ux636eux5e93ux8bbeux8ba1}

\subsubsection{4.5.1 数据库概念结构设计}\label{ux6570ux636eux5e93ux6982ux5ff5ux7ed3ux6784ux8bbeux8ba1}

系统数据库围绕用户、课程、任务、知识点、资源、评论消息、点赞关系和学习记录等核心实体展开。主要关系如图4-3所示。

\tikzDatabaseConceptFigure

从概念设计角度看，数据库结构并非简单地为页面展示提供存储空间，而是承担整个平台业务关系表达的职责。用户实体连接学习行为与互动行为，课程实体连接知识点、任务与资源，记录实体则承担学习过程留痕功能。这样的设计使系统不只能够保存结果数据，还能够保留过程数据，为后续的统计分析与教学干预提供基础。

\subsubsection{4.5.2 主要数据表设计}\label{ux4e3bux8981ux6570ux636eux8868ux8bbeux8ba1}

表4-1 主要数据表设计

数据库中的主要数据表如下。

\begin{longtable}[]{@{}ll@{}}
\toprule\noalign{}
数据表 & 说明 \\
\midrule\noalign{}
\endhead
\bottomrule\noalign{}
\endlastfoot
\texttt{users} & 存储用户基础信息、角色、学校、班级等 \\
\texttt{courses} & 存储课程信息和课程内容索引 \\
\texttt{contents} & 存储课程层级内容结构 \\
\texttt{tasks} & 存储学生任务代码、状态与时间信息 \\
\texttt{knowledges} & 存储知识点与内容关联信息 \\
\texttt{resources} & 存储教学资源与权限信息 \\
\texttt{comments} & 存储消息与评论内容 \\
\texttt{likes} & 存储点赞关系 \\
\texttt{records} & 存储学习行为和错误记录 \\
\end{longtable}

在逻辑设计上，\texttt{users} 表是权限识别与角色管理的基础；\texttt{courses}、\texttt{contents} 与 \texttt{knowledges} 共同构成课程内容组织体系；\texttt{tasks} 表用于承载学生任务完成状态与代码提交结果；\texttt{resources} 表用于管理教学材料；\texttt{comments} 与 \texttt{likes} 则用于维持互动关系；\texttt{records} 表负责记录页面访问、代码运行和错误类型等过程数据。各数据表之间通过主外键关系或业务关联字段建立联系，从而支撑课程学习、任务执行和行为分析等核心流程。

本人认为，数据库设计的关键不在于数据表数量多少，而在于业务语义是否清晰。若一个表承担过多职责，后续维护时就容易出现字段膨胀和语义混乱；若表之间关联关系过于松散，又会削弱数据利用价值。因此，本文在设计时尽量使每张表承担较明确的业务职责，并为后续扩展预留适当空间。

\section{智能辅学模块设计}\label{ux667aux80fdux8f85ux5b66ux6a21ux5757ux8bbeux8ba1}

智能辅学模块主要围绕``提问输入''\,``场景识别''\,``提示词构造''\,``模型生成''和``结果返回''五个环节展开。针对不同学习任务，系统将分别调用自由问答、代码解释、报错分析、问题分解、抽象建模和算法设计等功能接口，其处理流程如图4-4所示。

\tikzAiFlowFigure

为避免大模型输出与教学目标脱节，本文进一步将该模块细化为``场景识别---约束生成---教学化输出---风险兜底''的四层结构。

首先，在场景识别层，系统不将学生输入统一视为普通问答，而是结合问题文本、代码内容、运行结果和任务上下文，对请求进行类型划分。具体包括自由问答、代码解释、报错解析、问题分解、抽象建模和算法设计引导六类任务。这样做的目的，在于让不同任务调用不同的提示模板与输出结构，避免模型在所有场景下采用同一种回答方式。

其次，在约束生成层，系统将学生年级、课程主题、知识点标签、任务目标和回答边界写入提示内容。本人在设计中强调两类约束：一类是教学约束，即回答应当循序渐进、语言适龄、避免无前提地直接给出完整答案；另一类是安全约束，即拒绝危险代码、规避越权请求，并限制与课程目标无关的输出扩散。该思路与现有研究中关于人在回路和教学对齐的观点是一致的{[}3{]}{[}8{]}{[}11{]}。

再次，在教学化输出层，系统不追求一次性给出最完整答案，而是优先保证回答的教学可用性。以报错解析为例，系统要求模型依次给出错误含义、错误位置、可能原因与修改建议；以算法设计引导为例，系统要求模型先帮助学生明确输入、输出与关键步骤，再逐步组织流程。通过这种结构化输出，平台更容易把模型回答转化为学生能够理解和消化的学习材料。

最后，在风险兜底层，系统保留日志记录、人工复核与规则拦截机制。对于高风险输出、低置信度回答或涉及完整作业答案的请求，平台优先采用弱化回答、返回思路框架或提示教师介入的处理方式。本人认为，这一层并非附加功能，而是教育场景接入 LLM 的必要条件。缺少兜底机制的大模型，更像一辆制动不足的车，速度虽快，却不适合直接开进课堂。

\section{系统部署设计}\label{ux7cfbux7edfux90e8ux7f72ux8bbeux8ba1}

系统前端可通过 Node 环境进行 Nuxt 部署，后端可通过 Uvicorn 运行 FastAPI 服务。数据库使用 SQL Server，缓存使用 Redis，静态资源与教学视频使用阿里云 OSS 进行存储。后续若需要提升服务稳定性，可进一步接入负载均衡和弹性伸缩机制。

从部署关系看，前端应用主要负责页面渲染与交互请求发起，后端服务负责接口处理与业务调度，数据库负责结构化数据持久化，Redis 负责部分高频数据缓存，OSS 则负责教学视频等非结构化资源存储。该部署方式能够较清晰地划分职责边界，避免所有压力集中到单一服务节点。

在毕业设计的实现规模下，上述部署方式已经能够满足系统演示和功能联调需求；若面向更大规模的教学使用，则还需要进一步考虑服务容错、日志监控、备份恢复和并发调度等问题。本人认为，部署设计就像搭建一座教学楼的管线系统，平时不显眼，但一旦规划混乱，后续所有功能都会受到牵连。因此，在论文中明确部署关系，有助于完整说明系统从设计到运行的落地路径。

\chapter{系统详细实现}\label{ux7cfbux7edfux8be6ux7ec6ux5b9eux73b0}

\section{用户认证与权限控制实现}\label{ux7528ux6237ux8ba4ux8bc1ux4e0eux6743ux9650ux63a7ux5236ux5b9eux73b0}

系统采用 JWT 作为用户认证机制。前端在登录成功后保存用户状态信息，并通过全局认证中间件拦截需要登录权限的页面访问请求。后端通过依赖注入的方式完成用户身份校验，并对公开接口和受保护接口进行分组管理。

在实现思路上，认证模块的核心目标是保证``用户身份可信''和``接口访问有边界''。学生与教师虽然共享同一平台入口，但其可访问页面和可执行操作并不完全一致，因此系统在登录后不仅需要识别用户是否已认证，还需要识别其角色属性，并据此决定可见菜单、可访问路由和可调用接口范围。

前端侧主要负责维持登录态和页面级拦截。当用户完成登录后，系统保存必要的身份状态信息，并在页面切换时优先判断当前页面是否需要认证。这样做的好处在于，未经授权的访问会在前端阶段被提前拦截，从而减少无效请求进入后端。后端侧则承担最终校验职责，通过统一认证依赖校验令牌有效性和用户身份，并在关键接口中进一步判断角色权限。

本人认为，权限控制模块的价值不只是``谁能进来''，还在于``进来之后能做什么''。若缺少细粒度权限边界，教师端资源维护能力和学生端学习能力之间就可能发生混淆，进而影响数据安全与系统秩序。因此，认证与权限控制虽然在界面上不显眼，却是整个平台稳定运行的基础。

\section{课程与知识点管理功能实现}\label{ux8bfeux7a0bux4e0eux77e5ux8bc6ux70b9ux7ba1ux7406ux529fux80fdux5b9eux73b0}

平台通过课程分类选择器、课程搜索选择器和知识点选择器实现课程内容组织。教师可以在课程更新页面中完成课程内容编辑、知识点关联和代码内容维护，前端通过富文本编辑器提升课程内容录入效率，后端则负责课程与知识点数据的持久化处理。

课程与知识点管理模块的实现重点在于内容结构化。若课程内容只是零散堆放，学生在学习时就难以形成清晰的知识路径，教师在维护时也难以保证内容一致性。因此，平台将课程、知识点与资源之间建立关联关系，使课程不仅是展示页面，更是知识组织与任务承载的基本单位。

在教师使用流程中，课程维护通常包括新增课程、编辑课程描述、关联知识点、补充示例代码和调整内容顺序等操作。为了降低维护成本，前端页面需要尽量把这些操作集中在同一工作界面中，减少教师在多个页面间来回切换。后端则负责接收结构化表单数据，并保证课程与知识点之间的关联关系能够稳定写入数据库。

对于学生而言，该模块的价值主要体现在``看得懂课程结构''和``找得到对应内容''。本人认为，课程模块设计得是否合理，会直接影响后续任务管理、智能问答和资源推荐等模块的使用体验。课程结构清晰，后续功能才有稳定的落点；课程结构混乱，其他模块再强也容易沦为零散工具。

\section{在线编程与代码运行功能实现}\label{ux5728ux7ebfux7f16ux7a0bux4e0eux4ee3ux7801ux8fd0ux884cux529fux80fdux5b9eux73b0}

在线编程功能是系统核心模块之一。学生提交代码后，后端会先进行语法检查，再进行安全检查，检查通过后由安全执行器完成实际运行，并将执行结果或错误信息返回前端，同时将运行行为记录到数据库中。其主要流程如图5-1所示。

\tikzCodeRunSequenceFigure

从实现逻辑看，该模块可分为输入接收、静态检查、受控执行和结果回传四个阶段。输入接收阶段主要完成学生代码、任务标识和运行参数的统一收集；静态检查阶段主要识别语法错误、危险导入和高风险函数调用；受控执行阶段则负责在限定环境中运行代码，并监控运行时长与输出规模；结果回传阶段最终把标准输出、异常信息或检查失败原因反馈给前端页面。

其中，静态检查与受控执行是该模块的关键。若系统在接收代码后直接执行，平台稳定性和服务器安全性都将面临较大风险；若系统只给出``运行失败''而不提供清晰原因，又无法满足教学场景需要。因此，本平台在设计上同时追求两点：一是控制执行边界，二是保留教学所需的反馈信息。前者解决安全问题，后者解决学习问题，两者缺一不可。

此外，代码运行模块还承担行为记录功能。每一次运行请求，无论成功与否，都可能反映学生的学习状态。例如，多次重复出现同类语法错误，往往意味着该知识点尚未被真正掌握；某一任务长时间停留在运行失败状态，也可能意味着学生在任务分解阶段已经遇到困难。因此，代码运行结果不仅服务于当前页面反馈，也为教师后续分析学生学习过程提供了数据基础。

\section{大模型辅助教学功能实现}\label{ux5927ux6a21ux578bux8f85ux52a9ux6559ux5b66ux529fux80fdux5b9eux73b0}

\subsubsection{5.4.1 自由问答功能实现}\label{ux81eaux7531ux95eeux7b54ux529fux80fdux5b9eux73b0}

系统为编程相关问题提供流式问答能力。为了限制回答范围并贴合学生使用场景，后端在请求发送给模型之前增加了系统提示词，用于约束回答主题和回复方式。

在该功能中，系统并不把模型简单当作开放聊天接口，而是通过提示约束将其限定在编程学习支持场景内。也就是说，模型回答的重点应放在概念解释、思路提示和学习引导上，而不是生成与课程无关的泛化内容。这样做可以在一定程度上降低回答偏离教学目标的概率。

\subsubsection{5.4.2 代码解释功能实现}\label{ux4ee3ux7801ux89e3ux91caux529fux80fdux5b9eux73b0}

代码解释功能以学生提交的代码为输入，要求模型先概括程序整体功能，再逐步解释关键语句，从而帮助学生建立从整体到局部的理解过程。

本人在该功能设计中更强调解释顺序而非解释长度。对于初学者而言，若系统一开始就逐行解释所有细节，学生往往会陷入局部信息而失去整体把握。因此，本平台优先要求模型先说明程序``在做什么''，再说明``为什么这样做''，最后才进入关键语句解释。这种组织方式更符合中小学生由整体到局部的理解路径。

\subsubsection{5.4.3 报错解析功能实现}\label{ux62a5ux9519ux89e3ux6790ux529fux80fdux5b9eux73b0}

报错解析功能会同时接收学生代码与运行输出，并通过提示词约束模型按``三步式''输出，即错误含义解释、错误位置定位和启发式修改引导。这种方式更适合中小学生编程学习中的循序渐进式反馈。

与直接展示原始报错信息相比，报错解析功能更强调``翻译''和``拆解''。许多初学者并不是完全不能改错，而是读不懂错误信息，也不知道应先看哪一行、先改哪一步。因此，系统在此处引入结构化输出，不是为了增加模型参与感，而是为了把原本碎片化、技术化的报错信息改写为学生能够操作的修改线索。

\subsubsection{5.4.4 问题分解功能实现}\label{ux95eeux9898ux5206ux89e3ux529fux80fdux5b9eux73b0}

系统根据课程内容与学生任务描述，将复杂编程任务分解为若干可执行步骤，帮助学生降低问题理解难度。

问题分解功能主要服务于综合任务场景。当学生面对较长任务描述时，经常会出现``不知道从哪里开始''的情况。对此，系统不直接替学生完成程序，而是将任务拆解为输入分析、数据处理、条件判断、循环组织和结果输出等若干步骤，使学生能够先找到入口，再逐步推进实现过程。

\subsubsection{5.4.5 抽象建模功能实现}\label{ux62bdux8c61ux5efaux6a21ux529fux80fdux5b9eux73b0}

系统通过抽象建模功能引导学生从具体问题中提炼一般规律，帮助其建立更稳定的问题求解模式。

这一功能的重点不在于给出标准模型，而在于帮助学生建立``从具体到一般''的过渡。本人认为，中小学生在编程学习中最常见的问题之一，是能够在单个例题中模仿做法，却难以将方法迁移到相似问题。抽象建模功能的价值，正是在于通过提炼输入对象、处理规则和输出目标，帮助学生逐步形成可迁移的思维框架。

\subsubsection{5.4.6 算法设计引导功能实现}\label{ux7b97ux6cd5ux8bbeux8ba1ux5f15ux5bfcux529fux80fdux5b9eux73b0}

算法设计引导功能以``开始、输入、处理、循环、输出、结束''等模板提示学生组织思路，使其逐步形成结构化算法表达能力。

在实现上，该功能更接近思路支架而非答案生成器。系统通过固定的过程模板帮助学生组织解题步骤，使其先明确流程结构，再进入代码表达阶段。这样做的目的，是把原本容易混乱的解题过程变成可拆分、可检查的步骤链条，从而降低学生在编码前的思维混乱程度。

\section{师生互动与消息功能实现}\label{ux5e08ux751fux4e92ux52a8ux4e0eux6d88ux606fux529fux80fdux5b9eux73b0}

平台通过评论消息相关接口实现师生之间以及学生之间的互动。系统支持消息发送、消息读取、未读状态管理和对话内容展示，可用于课后答疑与学习交流。

该模块的实现意义在于补足教学平台中的``人际反馈链路''。如果平台只有课程、任务和模型回答，而缺少教师与学生之间的直接互动，那么许多需要情境判断和个别化解释的问题仍然难以得到有效处理。因此，消息模块并不是对智能功能的重复建设，而是对其必要补充。

在具体实现中，系统需要同时处理消息内容、发送关系、接收状态与未读标记等信息，使互动过程既可追踪又便于管理。对于教师而言，消息模块能够承载课后答疑与个别提醒；对于学生而言，则能够在模型回答不足或需要确认任务要求时，保留向教师继续求助的正式通道。

\section{教学资源管理功能实现}\label{ux6559ux5b66ux8d44ux6e90ux7ba1ux7406ux529fux80fdux5b9eux73b0}

教师端可对教学资源进行创建、修改、删除和权限设置。资源信息包括名称、内容、代码、教材页码、视频链接和知识点标签等，用于支撑课堂教学与课后扩展学习。

教学资源管理模块的重点在于资源与教学内容之间的可关联性。若资源只以文件形式独立存在，其教学价值往往取决于教师额外解释；若资源能够与课程、知识点和任务建立对应关系，学生在学习过程中就更容易找到``当前问题对应的辅助材料''。因此，本平台中的资源管理并非简单文件上传，而是面向教学组织的资源编排。

此外，资源权限设置也具有实际意义。部分资源适合全体学生查看，部分资源则更适合作为教师备课材料或分层教学补充内容。通过权限控制，平台能够在同一资源体系内兼顾公开学习支持与教师内部管理需求，从而提高资源复用效率。

\section{学习行为记录与事件追踪功能实现}\label{ux5b66ux4e60ux884cux4e3aux8bb0ux5f55ux4e0eux4e8bux4ef6ux8ffdux8e2aux529fux80fdux5b9eux73b0}

系统会记录用户学习过程中的页面访问、代码运行、错误类型等行为数据，并通过可视化页面展示学习记录与交互情况，为教师掌握学生学习进度和错误分布提供数据支撑。

从平台价值角度看，学习行为记录模块相当于系统的``观察窗''。若没有该模块，平台虽然能够提供教学功能，却很难回答``学生到底如何使用这些功能''这一关键问题。通过记录页面访问、代码运行、错误类型和交互过程，系统可以把原本不可见的学习过程转化为可分析的数据线索。

本人认为，该模块的意义不仅在于统计，更在于为后续教学干预提供依据。教师可以依据行为数据识别高频错误知识点，判断哪些任务难度偏高，或发现哪些学生长期缺乏有效互动。对毕业设计而言，这一模块已经为后续研究者继续引入学习者建模、分层推荐和教学效果评价提供了基础数据入口。

\chapter{系统测试与结果分析}\label{ux7cfbux7edfux6d4bux8bd5ux4e0eux7ed3ux679cux5206ux6790}

\section{测试环境}\label{ux6d4bux8bd5ux73afux5883}

表6-1 测试环境配置

\begin{longtable}[]{@{}ll@{}}
\toprule\noalign{}
类别 & 配置 \\
\midrule\noalign{}
\endhead
\bottomrule\noalign{}
\endlastfoot
前端环境 & Nuxt 3 + Node.js \\
后端环境 & FastAPI + Python \\
数据库 & SQL Server \\
缓存 & Redis \\
模型接口 & DashScope 等 \\
\end{longtable}

测试环境的设置应尽量与系统实际运行环境保持一致，以保证测试结果具有基本参考价值。前端环境主要负责页面渲染与交互逻辑验证，后端环境主要承担接口请求处理、权限校验和业务调度，数据库与缓存环境则用于验证数据读写和高频访问场景下的系统表现。模型接口环境主要用于测试智能辅学功能在真实调用链路中的可用性与响应完整性。

\section{功能测试}\label{ux529fux80fdux6d4bux8bd5}

后续可按下表补充功能测试结果。

表6-2 功能测试结果

\begin{longtable}[]{@{}lllll@{}}
\toprule\noalign{}
测试编号 & 测试模块 & 测试内容 & 预期结果 & 实际结果 \\
\midrule\noalign{}
\endhead
\bottomrule\noalign{}
\endlastfoot
T1 & 登录认证 & 正确账号密码登录 & 登录成功并跳转 & {[}待补充{]} \\
T2 & 在线编程 & 提交正确代码 & 返回正确输出 & {[}待补充{]} \\
T3 & 报错分析 & 提交错误代码 & 返回错误解释与定位 & {[}待补充{]} \\
T4 & 消息互动 & 发送一条消息 & 对方可正常接收 & {[}待补充{]} \\
T5 & 资源管理 & 新增教学资源 & 资源可保存并展示 & {[}待补充{]} \\
\end{longtable}

在测试组织方式上，功能测试应围绕核心业务链路展开，而不宜仅做零散页面点击检查。换言之，每一项测试不仅要验证单点功能是否可用，还应验证其前后流程是否闭环。例如，在线编程测试不仅要看代码是否能运行，还要同时检查运行记录是否写入、错误信息是否返回、前端页面是否正确展示结果。

此外，智能辅学相关测试应特别关注输出结构是否符合预期，而不仅是模型是否``返回了内容''。本人认为，教育场景中的模型测试更像改试卷，不仅要看有没有答题，还要看答题是否切题、是否有条理、是否适合当前学生理解。基于这一点，后续功能测试表中还可以增加对代码解释结构完整性、报错解析三步式输出完整性和问题分解层次性的检查项目。

\section{性能测试}\label{ux6027ux80fdux6d4bux8bd5}

系统性能测试可从接口响应时间、代码运行耗时、大模型响应时延和多用户并发访问等角度展开。考虑到本文当前阶段以系统设计与实现为主，性能测试更适合作为后续补充验证内容。后续可采用统一测试脚本对登录、课程访问、代码运行和智能问答等高频接口进行重复采样，并结合平均响应时间、峰值响应时间、错误率和并发吞吐量等指标进行统计分析。

\section{测试结果分析}\label{ux6d4bux8bd5ux7ed3ux679cux5206ux6790}

从当前系统联调情况来看，平台主要功能链路已经能够完成闭环运行，说明整体架构设计与模块划分具备基本可行性。对于登录认证、课程访问、在线编程、报错解析、消息互动和资源管理等核心功能，只要后续将真实测试结果补入表格，即可进一步形成较为完整的功能验证证据。

需要说明的是，本文现阶段尚未形成成体系的课堂实验数据和大规模性能测试数据，因此本章分析应严格限定在系统实现与功能可用性层面，而不宜直接推导平台对学生学习效果的最终影响。后续若开展真实教学实验，可结合学生任务完成率、错误修正效率、学习兴趣量表和自我效能感量表等指标，对平台的教学支持效果进行更具说服力的评价{[}4{]}。

\chapter{结论}\label{ux7ed3ux8bba}

本文围绕中小学生编程学习场景，完成了一套基于大模型的智能编程学习平台的设计与实现。系统采用前后端分离架构，围绕课程管理、任务学习、在线编程、智能辅学、消息互动、资源管理和行为记录等模块展开建设，并重点对智能辅学模块的场景划分、提示约束和输出组织进行了面向教学目标的改造。

与将大模型简单接入聊天接口的实现方式相比，本文更强调教学场景适配、系统工程约束与教师可控性。本人认为，面向中小学生的编程辅学平台，其关键不在于模型是否足够强，而在于平台是否能够把模型能力转化为可理解、可追踪、可干预的教学支持过程。基于这一思路，本文的重点并非追求模型替代教师，而是构建教师、学生与大模型协同参与的学习支持机制。

需要说明的是，本文当前能够确认的结论主要集中于系统设计合理性与核心功能实现层面。至于平台对学生编程能力、学习兴趣、自我效能感及长期学习行为的实际影响，仍需依赖后续真实课堂实验、对照测试与过程数据分析进一步验证。后续研究可继续围绕课程知识库检索增强、学习者建模、提示策略评估、教师干预策略和教学实验设计等方向展开。

\chapter*{参考文献}\addcontentsline{toc}{chapter}{参考文献}\label{ux53c2ux8003ux6587ux732e}

{[}1{]} Wang S, Xu T, Li H, et al.~Large Language Models for Education: A Survey and Outlook{[}EB/OL{]}. arXiv, 2024{[}2026-04-05{]}. https://doi.org/10.48550/arXiv.2403.18105.

{[}2{]} Raihan N, Siddiq M L, Santos J C S, et al.~Large Language Models in Computer Science Education: A Systematic Literature Review{[}C{]}//Proceedings of the 56th ACM Technical Symposium on Computer Science Education V. 1. New York: ACM, 2025: 938-944. https://doi.org/10.1145/3641554.3701863.

{[}3{]} Pitts G, Hridi A P, Lekshmi-Narayanan A B. A Survey of LLM-Based Applications in Programming Education: Balancing Automation and Human Oversight{[}EB/OL{]}. arXiv, 2025{[}2026-04-05{]}. https://doi.org/10.48550/arXiv.2510.03719.

{[}4{]} Tang B, Liang J, Hu W, et al.~Enhancing Programming Performance, Learning Interest, and Self-Efficacy: The Role of Large Language Models in Middle School Education{[}J{]}. Systems, 2025, 13(7): 555. https://doi.org/10.3390/systems13070555.

{[}5{]} Chen L, Xiao S, Chen Y, et al.~ChatScratch: An AI-Augmented System Toward Autonomous Visual Programming Learning for Children Aged 6-12{[}C{]}//Proceedings of the CHI Conference on Human Factors in Computing Systems. New York: ACM, 2024: 1-19. https://doi.org/10.1145/3613904.3642229.

{[}6{]} Ma Q, Shen H, Koedinger K R, et al.~How to Teach Programming in the AI Era? Using LLMs as a Teachable Agent for Debugging{[}C{]}//Artificial Intelligence in Education. Lecture Notes in Computer Science. Cham: Springer Nature Switzerland, 2024: 265-279. https://doi.org/10.1007/978-3-031-64302-6\_19.

{[}7{]} Gonzalez-Maldonado D, Liu J, Franklin D. Evaluating GPT for use in K-12 Block Based CS Instruction Using a Transpiler and Prompt Engineering{[}C{]}//Proceedings of the 56th ACM Technical Symposium on Computer Science Education V. 1. New York: ACM, 2025: 388-394. https://doi.org/10.1145/3641554.3701910.

{[}8{]} Pirzado F A, Ahmed A, Mendoza-Urdiales R A, et al.~Navigating the Pitfalls: Analyzing the Behavior of LLMs as a Coding Assistant for Computer Science Students---A Systematic Review of the Literature{[}J{]}. IEEE Access, 2024, 12: 112605-112625. https://doi.org/10.1109/ACCESS.2024.3443621.

{[}9{]} Humble N. Risk management strategy for generative AI in computing education: how to handle the strengths, weaknesses, opportunities, and threats?{[}J{]}. International Journal of Educational Technology in Higher Education, 2024, 21(1). https://doi.org/10.1186/s41239-024-00494-x.

{[}10{]} Tlili A, Shehata B, Adarkwah M A, et al.~What if the devil is my guardian angel: ChatGPT as a case study of using chatbots in education{[}J{]}. Smart Learning Environments, 2023, 10(1). https://doi.org/10.1186/s40561-023-00237-x.

{[}11{]} Jeon J B, Lee S. Large language models in education: A focus on the complementary relationship between human teachers and ChatGPT{[}J{]}. Education and Information Technologies, 2023, 28(12): 15873-15892. https://doi.org/10.1007/s10639-023-11834-1.

{[}12{]} Yim I H Y, Su J. Artificial intelligence (AI) learning tools in K-12 education: A scoping review{[}J{]}. Journal of Computers in Education, 2024. https://doi.org/10.1007/s40692-023-00304-9.

\chapter*{致谢}\addcontentsline{toc}{chapter}{致谢}\label{ux81f4ux8c22}

{[}此处补充{]}

\appendix\chapter{附录}\label{ux9644ux5f55}

\section{附录 A 主要接口说明}\label{ux9644ux5f55-a-ux4e3bux8981ux63a5ux53e3ux8bf4ux660e}

{[}此处补充{]}

\section{附录 B 主要数据表字段说明}\label{ux9644ux5f55-b-ux4e3bux8981ux6570ux636eux8868ux5b57ux6bb5ux8bf4ux660e}

{[}此处补充{]}

\section{附录 C 关键代码片段}\label{ux9644ux5f55-c-ux5173ux952eux4ee3ux7801ux7247ux6bb5}

{[}此处补充{]}
